\section{Feedback and Proposed Novelty}
\label{sec:feedback_novelty}

\paragraph{Feedback Since Presentation 1}
% Action items and feedback received from the course instructors and teaching assistants:
[Placeholder: Insert specific feedback received following Presentation 1, including requested revisions regarding experimental scope, baseline coverage, and presentation structure.]

\paragraph{Weekly Mentor Progress}
% Summary of weekly progress meetings, scope adjustments, and milestone tracking:
[Placeholder: Insert summary of weekly mentor interactions, scope refinement milestones, and verified empirical progress reports.]

\paragraph{Reproduction Scope and Future Novelties}
In this investigation, we have not introduced novel architectural parameterizations or newly designed loss functions. Rather, our primary objective is an exhaustive, rigorous reproduction and empirical cross-paradigm benchmark of established post-processing transformations from the representation learning literature \citep{mu2018all,timkey2021all,su2021whitening,diera2024isotropy,ren2025correcting,li2026spectral}. 

We evaluate these historical methods systematically across newer sub-1B parameter models released in late 2024 and 2025, modern MTEB v2 benchmarks, and across two distinct operational tiers: dedicated contrastive embedding encoders and un-fine-tuned base causal autoregressive language models. 

Building upon the empirical baselines and diagnostic insights established in this comprehensive reproduction, our future work will explore hybridizing coordinate-wise standardization, mean-subspace projection deflation, and signal-to-noise ratio spectral tempering. By unifying these complementary techniques into an integrated pipeline, we plan to formulate novel training-free post-processing operators that expand the Pareto frontier of retrieval accuracy and dimensional compression beyond individual precursor methods.
